% answers_short.tex —— 2005 年·普通高等学校招生考试·湖北卷·理科 · 简短答案速查
% 由 tools/make_short_answers.py 自 answers.tex 的 \gksolhead/\gkanswer 提取，
% 并与 2005/2005_数学真题_湖北理科_答案要点.md 逐题比对后生成，勿手改（改源 MD 后重跑）。

\begin{gknotice}
\textbf{使用说明}：本节只给最终答案便于快速核对，完整推导见解析版。答案为编写者独立推导，\textbf{非官方标准答案}。
\end{gknotice}

\gkpart{选择题}
\begin{center}\begin{tabular}{ccccccccccccc}
\toprule
题号 & 1 & 2 & 3 & 4 & 5 & 6 & 7 & 8 & 9 & 10 & 11 & 12 \\
\midrule
答案 & B & B & D & C & A & B & C & C & D & C & D & A \\
\bottomrule
\end{tabular}\end{center}

\gkpart{填空题}
\begin{itemize}[leftmargin=2.6em]
  \item[\textbf{13.}] $[-6, 2]$。
  \item[\textbf{14.}] $\dfrac{63\sqrt{2}}{2}$。
  \item[\textbf{15.}] $-2$。
  \item[\textbf{16.}] $500$。
\end{itemize}

\gkpart{解答题}
\begin{itemize}[leftmargin=2.6em]
  \item[\textbf{17.}] $[5, +\infty)$。
  \item[\textbf{18.}] $\dfrac{\sqrt{70}}{14}$。
  \item[\textbf{19.}] $\xi$ 的分布列为 $P(\xi=1)=0.6, P(\xi=2)=0.28, P(\xi=3)=0.096, P(\xi=4)=0.024$；数学期望 $\mathbf{E}\xi = 1.544$；李明在一年内领到驾照的概率为 $0.9976$。
  \item[\textbf{20.}] (1) 直线 $AC$ 与 $PB$ 所成角的余弦值为 $\dfrac{3\sqrt{7}}{14}$；(2) $N$ 点到 $AB$ 的距离为 $1$，$N$ 点到 $AP$ 的距离为 $\dfrac{\sqrt{3}}{6}$
  \item[\textbf{21.}] (1) $\lambda$ 的取值范围是 $(12, +\infty)$；直线 $AB$ 的方程为 $x + y - 4 = 0$；(2) 存在，四点 $A, B, C, D$ 对任意 $\lambda \in (12, +\infty)$ 恒共圆，证明见解析
  \item[\textbf{22.}] (1) 证明见解析；(2) 数列 $\{a_n\}$ 存在极限，且 $\lim_{n\to\infty} a_n = 0$；(3) 取正整数 $N = 1023$（或取满足 $N \geqslant 1023$ 的任意正整数），当 $n > N$ 时，对任意 $b > 0$ 恒有 $a_n < \dfrac{1}{5}$
\end{itemize}
